\chapter{L’eterno mutamento}

«A cosa stai pensando?»\\
«Niente. Guardavo le stelle.»\\
«Che costellazione vedi, Giovanni?»\\
«Non le riconosco pi\`u. Pensi che i cieli stiano cambiando?»





\par\medskip
Gli occhi di Giovanni sono chiusi e la testa \`e rivolta al cielo.
Il Grande Carro \`e quasi sdraiato sotto la Stella Polare,
lui invece \`e sdraiato sulla rete del vagabondo.
Fa caldo, ma hanno acceso un fuoco dentro uno dei bidoni.

\par\medskip
«Quando mai sono rimasti uguali a se stessi?»\\
«E ora? In cosa muteranno?»\\
«Te lo chiedi perch\'e il mondo \`e tornato a interessarti?
Vuoi capire se tu puoi averci ancora a che fare?»\\
«Quante domande fai per essere uno che non ama parlare?
Non lo so. \`E che mi chiedo se questa volta il mondo sia cambiato pi\`u del solito.
\`E lo stesso cambiamento o \`e un cambiamento diverso?»\\
«Tu cosa pensi? So che se me lo chiedi...»\\
«No, non ho gi\`a una mia idea. Ho rinunciato alle idee.
Le idee mutano in fretta. \`E inutile possederle.
Chiamiamo idea un processo mentale...»\\
«Già, ma dimmi, Giovanni, quale mutamento delle tue abitudini ti ha cos\`i turbato?
Senti bisogno di una \emph{cate...}chesi?»\\
«Lascia perdere. Quando la vista ha iniziato a calarmi \`e stato l’olfatto il senso che ho sviluppato maggiormente.
Un odore pu\`o parlarmi per ore e io posso stare ore ad ascoltarlo.»

\par\medskip
«Ad ascoltarla, vuoi dire. Io comunque odorerei volentieri una bella torta.
Il dolce \`e l’unica cosa che mancava stasera.»

\par\medskip
La Stella Polare \`e rimasta ferma, mentre le sette stelle del Carro si trovano in una posizione diversa ora.
La loro distanza relativa \`e la stessa: nel complesso l’intera costellazione non ha cambiato forma.
\`E stata una trasformazione isometrica, una rotazione delle stelle attorno alla Polare.
Per questo lei \`e rimasta ferma.

\par\medskip
«Comunque io credo che il mondo, la biosfera, Gaia, gli d\`ei, il Tao, tutto si trasformi,
ma che qualcosa rimanga sempre tale e quale.
E quel qualcosa sai cos’\`e?»

\par\medskip
Giovanni ascolta l’amico e si gira verso di lui.
Questa piega della conversazione deve averlo incuriosito.

\par\medskip
«Il tempo?»

\par\medskip
\begin{center}
\includegraphics[width=1.0\linewidth]{planisfero senza sfondo.png}
\end{center}
